% 临时英文行文稿的译文：进一步的结构性结果与失效的归约。
% 来源标识见 provenance/SOURCE_MAP.json，R15,R21,R26,R27。

\section{为什么有限个点不能代替一根完整针}\label{sec:finite-samples}
在优化模型中要求整个区间都属于该体，代价很高。只检查端点、中点，或一份预先固定的细密内部点列表，看起来是可控的简化。下面的构造说明，即使不离散化方向变量，这种简化也会丢掉什么。

\begin{theorem}[有限采样的能力上限]\label{thm:sample-ceiling}
固定有限集 $T\subset[0,1]$、高度 $0<h_0<1/4$ 和 $\varepsilon>0$。存在紧的原点星形集 $K\subset B_1$，它是有限个闭圆扇形的并，满足 $|K|<\varepsilon$，而且对每个有向方向 $q$，都存在 $s\in(-1/4,1/4)$ 使得
\[
 \operatorname{Rot}_q(s+t-1/2,h_0)\in K\quad\text{对所有 }t\in T.
\]
可以用有限个闭方向区间覆盖整个圆周，并在每个区间上把 $s$ 取为常数，同时保留所有边界记录。各处都使用同一个预先指定的正高度。
\end{theorem}

\begin{proof}
必要时把 $0,1$ 加入 $T$。在物理角圆周上选一个开稠密子集 $A$，使 $|A|<\min(2\varepsilon,\pi)$。具体地，用开弧包围一个可数稠密集，并使这些开弧的总长度小于上述界。

对固定的 $q$ 和 $t$，考虑
\[
 f_{q,t}(s)=q+\arg(s+t-1/2+ih_0),\qquad -1/4<s<1/4.
\]
它的导数为
\[
 f'_{q,t}(s)=-\frac{h_0}{(s+t-1/2)^2+h_0^2}<0.
\]
因此，每个非空开子区间的像都包含一段开弧。开稠密集 $A$ 的原像是开稠密的。将有限集 $T$ 对应的这些原像取交，再从交集中选取 $s_q$。

固定 $s_q$，让 $q$ 变化。在原方向的某个开邻域 $V_q$ 内，有限个采样点的角度仍然落在 $A$ 中。由紧致性并收缩开覆盖，可得到有限个覆盖整个圆周的闭圆周区间 $I_j$，每个区间都包含在一个这样的邻域中。给 $I_j$ 配上相应的数值 $s_j$。定义
\[
 K=\bigcup_j\bigcup_{t\in T}
 \{r\operatorname{Rot}_q(s_j+t-1/2,h_0):0\le r\le1,\ q\in I_j\}.
\]
并中的每一项都是一个闭圆扇形。每个非零点的角度都在 $A$ 中，且其半径小于一，因为 $|s_j+t-1/2|<3/4$ 且 $h_0<1/4$。所以 $K$ 是紧的原点星形集，面积至多为 $|A|/2<\varepsilon$。对每个方向，总有一个包含它的闭区间提供所需的记录，接缝处也包括在内。
\end{proof}

两个采样端点的距离恰好为一。垂足位于线段内部，高度保持固定且为正。然而，当 $\varepsilon<h_0/2$ 时，上述候选记录的完整线段没有一条包含在 $K$ 中：否则星形性会填满一个面积为 $h_0/2$ 的三角形。这个反证只排除指定高度 $h_0$ 的那些候选记录；其他可能的单位针不必有这个高度，不能一并排除。障碍来自用有限列表替换所有内部点，而不是来自遗漏方向或让高度趋于零。

采样列表变化时，集合 $K$ 也随之变化。对一个固定的紧集 $K$，如果同一条记录在 $[0,1]$ 的一个稠密子集上对应的点都属于 $K$，那么由连续性和闭性，整条记录都属于 $K$。因此，该定理仍为精确的连续统论证，或带有已证明的一致误差界的自适应逼近，留下了余地。它排除的是固定有限采样模型本身的正面积下界定理。

\section{替换远处的针时，可以显式控制半径误差}\label{sec:radius}
研究下确界时，半径无界是另一个看似难以处理的障碍。这里可以作可控的归约，但记录数量仍不受控制。

\begin{proposition}[半径截断]\label{prop:radius}
设 $K$ 是紧的原点星形集，在每个射影方向都包含一条完整的单位线段，记 $A=|K|\le U$。若 $R\ge4U+2$，则存在在所有方向都包含单位线段的有界 Borel 星形集 $F\subset\overline B_R$，使得
\[
 |F|\le A(1+340/R).
\]
从而，若 $A_R$ 是 $\overline B_R$ 中这类 Borel 体的面积下确界，则
\[
 A_*\le A_R\le(1+340/R)A_*.
\]
\end{proposition}

\begin{proof}
记 $\mathrm{Bad}_R$ 为那些在 $K\cap\overline B_R$ 中没有完整单位线段的方向。由单位线段端点的紧致性，可行方向集是闭的，所以 $\mathrm{Bad}_R$ 是开的。保留 $K\cap\overline B_R$，并在每个坏方向上加入以原点为中点的单位线段。所加线段的并的面积为 $|\mathrm{Bad}_R|/4$：每个射影方向贡献两条长度为 $1/2$ 的物理射线。因此，只需估计坏方向的测度。

记 $D_s$ 为这样的射影方向集：对应的物理射线中，至少有一条在 $K$ 内的径向延伸长度达到 $s$。极坐标积分给出 $|D_s|\le2A/s^2$。在一个坏方向上，任取 $K$ 中实际存在的一条单位线段，并在沿该方向定向的坐标系中，把端点写成 $(c-1/2,h),(c+1/2,h)$，其中 $c\ge0$。它有一点落在 $\overline B_R$ 外，所以其中点的范数大于 $R-1/2$。它与原点张成的三角形给出 $|h|\le2A\le2U$。对 $R$ 的假设蕴含 $c\ge R/2\ge1$ 且 $|h|\le c$。

当 $h=0$ 时，该方向本身就在 $D_{R/4}$ 中。对 $h>0$，整个底边张出的角区间为
\[
 I=\left[\arctan\frac{h}{c+1/2},\arctan\frac{h}{c-1/2}\right]
 \subset D_{c/2}.
\]
最后的包含关系用到了整个三角形：每条中间射线都在至少为 $c-1/2\ge c/2$ 的半径处到达底边。令 $J$ 从单位线段的方向延伸至 $I$ 的远端。对角度导数积分得到
\[
 |I|\ge\frac{h}{(c+1/2)^2+h^2},\qquad
 |J|\le\frac{h}{c-1/2},\qquad
 \frac{|J|}{|I|}\le\frac{13c}{2}.
\]
高度为负时，对同一个物理体作反射后的角度论证即可。

令 $C_k=2^kR/2$。若 $C_k\le c<2C_k$，则 $I\subset D_{C_k/2}$ 且 $|J|/|I|<14C_k$。因此，该坏方向属于
\[
 \left\{\mathcal M_{\mathrm{RP}}1_{D_{C_k/2}}>\frac1{14C_k}\right\},
\]
其中 $\mathcal M_{\mathrm{RP}}$ 是射影圆周上通常的非中心区间平均极大算子。由于不等式是严格的，若该方向恰为端点，就可以把 $J$ 稍微扩大。

应用式~\eqref{eq:circle-maximal-weak} 的圆周弱型估计，这一尺度贡献的方向测度至多为 $336A/C_k$。将各尺度的界相加，再加上零高度方向，得到
\[
 |\mathrm{Bad}_R|\le\frac{1344A}{R}+\frac{32A}{R^2}.
\]
将中心补针的并记为 $C_R$。其面积与实际新增量满足
\[
 |C_R\setminus(K\cap\overline B_R)|\le |C_R|
   =\frac{|\mathrm{Bad}_R|}{4}
   \le\frac14\left(\frac{1344A}{R}+\frac{32A}{R^2}\right)
   \le\frac{340A}{R},
\]
其中最后一步使用 $R\ge2$。净面积变化还要减去被截掉的旧材料 $|K\setminus\overline B_R|$；因此上述新增量的上界足以推出 $|F|\le A(1+340/R)$，第一项断言得证。将它用于逼近 $A_*$ 且面积小于 $U$ 的原始有限体，得到 $A_R\le(1+340/R)A_*$；对任意容许的 Borel 体应用引理~\ref{lem:finite-open-recovery}，则得 $A_*\le A_R$。
\end{proof}

截断后，可以在 $B_{R+1}$ 内作有限多边形或有理多边形恢复；额外的半径提供了开放余量。因此，若对恢复类的\emph{所有复杂度}都有一致下界 $c_{R+1}$，就能推出
\[
 A_*\ge\frac{c_{R+1}}{1+340/R}.
\]
只检查固定数量的胞腔而得到的下界不满足这个假设。截断估计将向远处逃逸与复杂度分开处理，并未同时解决这两个问题。

\section{完备的强面积空间与一列有用的近极小体}\label{sec:strong-metric}
要系统地使用替换，可以先选取具有全局比较性质的近极小体。径向剖面的强收敛几乎已经提供了所需的框架。剩下的是一个几何问题：完整单位线段，尤其是那些高度塌缩到零的线段，必须在极限面积等价类的某个实际代表元中保留下来。

令 $\mathcal X$ 由有限面积 Borel 原点星形集的等价类组成，其中以平面对称差为零测集定义等价，并要求每个类都有一个 Borel 代表元，在每个射影方向上包含一条完整的闭单位线段。代表元可以无界。置
\[
 d(F,G)=|F\triangle G|,\qquad a(F)=|F|.
\]
关于代表元的要求是定义的一部分。单位线段上几乎处处属于该体，不能代替整个闭区间都包含在该体中。

\begin{theorem}[强面积闭性与不受限制的交换]\label{thm:strong-exchange}
度量空间 $(\mathcal X,d)$ 是完备的，且 $\inf_{\mathcal X}a=A_*$。对每个 $\alpha>0$ 及满足 $a(F)\le A_*+\alpha^2$ 的 $F\in\mathcal X$，都存在 $G\in\mathcal X$ 使得
\[
 a(G)\le a(F),\qquad d(G,F)\le\alpha,
 \qquad a(G)\le a(Q)+\alpha d(G,Q)\quad(Q\in\mathcal X).
\]
等价地，同一个 $G$ 满足
\begin{equation}\label{eq:relative-exchange}
 (1-\alpha)|G\setminus Q|\le(1+\alpha)|Q\setminus G|
 \quad\text{对每个 }Q\in\mathcal X.
\end{equation}
\end{theorem}

\begin{proof}
为一个 Cauchy 序列 $F_n$ 选取实际的全方向 Borel 代表元。利用有理半径处的成员关系取上确界，可知它们的径向延伸长度 $\rho_n$ 是 Borel 的。将每个径向区间闭合，这只改变一个面积为零的径向边界图像，并保留原来的每条单位线段。总面积有限，因此无限延伸的射线所占的角度构成一个 Borel 零测集。在这些零测集之外令 $Y_n=\rho_n^2$，在其上令它等于零。极坐标积分给出
\[
 d(F_n,F_m)=\frac12\int_0^{2\pi}|Y_n-Y_m|\,d\theta.
\]
其 $L^1$ 极限可选取一个非负、有限的 Borel 版本 $Y$。它的闭径向体 $S=\{0\}\cup\{re_\theta:0<r\le\sqrt{Y(\theta)}\}$ 是 $F_n$ 的面积极限。固定 $M>\sup_n|F_n|$。

取一个子序列，使通常的极大误差 $\mathcal M(|Y_n-Y|)$ 几乎处处趋于零。弱型极大估计配合一个 $L^1$ 误差可求和的子序列，保证可以这样选取。对两个物理提升，去掉所有无限射线集，以及 $Y$ 和可数个 $Y_n$ 的所有 Lebesgue 值例外集。剩下的射影方向集 $H$ 是 Borel 集，且具有全测度。

所选单位线段没有共同的半径界。我们只需一个较弱的事实：几乎每个方向都有无穷多个落在某个固定有界区域内的单位线段见证。对每个 $n$ 和 $R\ge4M+2$，由 $F_n$ 的实际径向延伸长度定义 $D_{s,n}$，其中也包括无限延伸的情形，并置
\[
 \mathcal B_{R,n}=D_{R/4,n}\ \cup\!\bigcup_{k\ge0}
 \left\{\mathcal M_{\mathrm{RP}}1_{D_{C_k/2,n}}>\frac1{14C_k}\right\},
 \qquad C_k=2^kR/2.
\]
这个显式定义的集是 Borel 的。命题~\ref{prop:radius} 中关于远针的证明，可逐点应用于 Borel 星形集中实际存在的任意单位线段：它与原点张成的三角形给出 $|h|\le2|F_n|$，其完整底边给出同样的角区间。这些几何估计没有使用紧致性。因此，若 $F_n\cap\overline B_R$ 在方向 $q$ 上不包含完整单位线段，则 $q\in\mathcal B_{R,n}$，并且
\[
 |\mathcal B_{R,n}|\le1344M/R+32M/R^2.
\]
对任意 Borel 体，我们不声称其无单位线段方向集可测；测量的是上面显式构造的 Borel 包含集。

选取 $R_\ell\to\infty$，并令 $Z=\bigcap_\ell\liminf_n\mathcal B_{R_\ell,n}$。由 Fatou 引理，每个下极限集的测度都不超过上式在 $R_\ell$ 处的量，所以 $Z$ 是零测集。对 $q\notin Z$，存在某个固定的 $R_\ell$，使无穷多个 $n$ 满足 $q\notin\mathcal B_{R_\ell,n}$。沿这些指标，在 $F_n\cap\overline B_{R_\ell}$ 内都有方向为 $q$ 的完整单位线段见证。它们的中点有一个收敛子序列。半径和子序列都可以依赖于 $q$；下面固定一个方向来论证，这就足够了。

固定 $q\in H\setminus Z$ 及这样一个收敛的单位线段序列。如果其极限高度非零，那么这些线段与原点张成的三角形按对称差面积收敛到一个非退化三角形 $T$。由强面积收敛，得到 $|T\setminus S|=0$。

现在设极限高度为零，并将其纵向区间写成 $[b-1/2,b+1/2]$。在该区间的内部取一个正数 $x_0$，再取 $\delta>0$，使 $[x_0-\delta,x_0+\delta]$ 仍落在区间的正半轴部分内。对足够大的 $n$，这整个纵向区间都属于第 $n$ 条单位线段。当 $h_n\ne0$ 时，其物理角像 $I_n$ 是该区间在映射 $x\mapsto q+\arctan(h_n/x)$ 下的像。对 $I_n$ 中几乎每个角度，都有 $Y_n\ge(x_0-\delta)^2$；有限值版本 $Y_n$ 可能在无限射线的零测集上不同于实际径向延伸长度的平方。将 $I_n$ 延长为包含 $q$ 的 $J_n$。导数 $-h_n/(x^2+h_n^2)$ 表明，对这些固定的 $x_0,\delta$，比值 $|J_n|/|I_n|$ 保持有界，而两个区间都收缩到 $q$。因此
\[
 \operatorname{avg}_{I_n}|Y_n-Y(q)|
 \le C\left(\mathcal M(|Y_n-Y|)(q)
                 +\operatorname{avg}_{J_n}|Y-Y(q)|\right)\longrightarrow0.
\]
好点条件推出 $Y(q)\ge(x_0-\delta)^2$。让 $\delta$ 递减，并让 $x_0$ 趋于正端点。对纵向区间的负半轴部分，使用相反的物理提升，得到负端点的界。如果沿一个子序列有 $h_n=0$，则在好的 Lebesgue 点处，由 $Y_n(q)\to Y(q)$ 和实际存在的径向单位线段直接得到同一结论。因此，整个极限零高度单位线段都在 $S$ 中，即使它的垂足位于底边之外也如此。

最后，我们同时恢复所有非退化单位线段，而不取 $S$ 的平面闭包。令 $E$ 为所有满足 $|T\setminus S|=0$ 的非退化三角形 $T$ 的并，其中每个三角形都由原点与一条单位线段张成。把它们的位姿限制为端点范数不超过某个整数 $L$、高度绝对值至少为 $1/m$。由于三角形的对称差面积连续地依赖于端点，面积亏损为零的位姿集是紧的。这些位姿对应的三角形的并也是紧的。对 $L,m$ 取可数并，就可知 $E$ 是 Borel 集。

令 $O$ 为如下点组成的开集：该点有一个球邻域 $B$，满足 $|B\setminus S|=0$。用可数基可证 $|O\setminus S|=0$。定义 $E$ 的每个三角形的内部都在 $O$ 中。在非退化三角形的边或顶点处，一个具有固定正开角的局部锥仍然进入其内部；这样的点不是 $O$ 的补集的密度为一的点。由密度定理得到 $|E\setminus O|=0$，因而 $|E\setminus S|=0$。这样便通过同一个开集 $O$ 处理了不可数多个三角形边界的并。

$H\setminus Z$ 之外的例外射影方向集是 Borel 零测集。在两个物理提升上加入这些方向的、以原点为中点的完整单位线段。它们的并 $C$ 是 Borel 原点星形集，面积为零。代表元 $S\cup E\cup C$ 在每个方向都包含一条完整单位线段，且与 $S$ 只相差一个零测集。原来的 Cauchy 序列收敛到这个代表元所在的等价类，完备性得证。

每个原始有限体都属于 $\mathcal X$；反过来，对每个类选取定义所要求的全方向 Borel 代表元，引理~\ref{lem:finite-open-recovery} 就给出任意小向上面积误差的原始有限体。因此，两个下确界相同。面积泛函非负，且关于 $d$ 是 $1$-Lipschitz 的。应用 Ekeland 变分原理~\cite{Ekeland1974}，取目标误差为 $\alpha^2$、允许距离为 $\alpha$，便得到断言中的 $G$ 及它与每个 $Q$ 的比较。代入“加入面积减去撤去面积”的精确恒等式，即得~\eqref{eq:relative-exchange}。
\end{proof}

对 $0<\alpha<1$，固定这个 $G$ 的一个实际 Borel 代表元。在 $G$ 上赋权 $1-\alpha$、在 $G$ 外赋权 $1+\alpha$，它在整个 $\mathcal X$ 上极小化这一加权面积，因为
\[
 \int_Q w_G-\int_G w_G
   =(1+\alpha)|Q\setminus G|-(1-\alpha)|G\setminus Q|\ge0.
\]
这个正权重一经固定，就不随竞争体改变。依次选取 $\alpha\to0$ 和满足定理误差要求的近极小输入，所获 $G$ 的面积趋于原问题的同一个下确界 $A_*$，且每个 $G$ 都满足相应的全局相对交换律。完备性并不保证这个序列有收敛子序列，加权结论也不保证存在无权极小体。余下的尖锐下界必须把交换律与实际的完整针几何结合起来。

\section{下确界能否达到，是另一个问题}\label{sec:attainment}
一种看似诱人的简化，是要求近极小体在某个方向的最长弦恰好长一。未经修改的近极小序列不一定满足这个要求。略微膨胀既能保持面积接近极小值，又能让每个方向都有更长的弦。膨胀后的体仍然属于原来的有限语法：对 $1<\lambda<2$，用高度为 $\lambda h$、中心为 $\lambda b\pm(\lambda-1)/2$ 的两条单位记录，替换一条膨胀后的源记录。它们相互重叠的底边区间覆盖长度为 $\lambda$ 的整个膨胀底边，所以所得星形集恰好是 $\lambda K$。

沿用第~\ref{sec:chords} 节的最长完整弦函数 $L_K$，对一个原始有限体令 $L_*(q)=\lim_{\delta\downarrow0}\inf\{L_K(p):\operatorname{dist}(p,q)<\delta\}$，并令 $Z=\{L_*=1\}$。这里必须取下半连续包络：孤立的长弦不能保证邻近方向也有长度余量。先设 $Z$ 非空且不是整个圆周。选取一个有界半代数选择器，逐方向选出实际单位线段；令 $\Gamma$ 为这个选择器图像的闭包，保留每个单侧端点极限。定义紧的原点星形集
\[
 H_\delta=\bigcup_{(q,N)\in\Gamma,\ \operatorname{dist}(q,Z)\le\delta}
                       \operatorname{conv}(0,N),\qquad
 H=\bigcup_{(q,N)\in\Gamma,\ q\in Z}\operatorname{conv}(0,N).
\]
它们都包含在 $K$ 中，且当 $\delta\downarrow0$ 时有 $H_\delta\downarrow H$。由方向、端点及凸组合参数的紧致性，既可证明这个极限等式，也可证明 Hausdorff 收敛。

若 $|H|<|K|$，则极坐标积分给出一个正测度的物理角集合 $B$ 以及 $a>0$，使得对每个 $\phi\in B$，原来的径向端点 $\rho_K(\phi)u_\phi$ 到 $H$ 的距离至少为 $a$，并且 $\rho_K(\phi)\ge a$。固定 $K\subset B_R$，选取 $\delta$ 使 $H_\delta\subset H+B_{a/4}$，再利用下半连续性，在距 $Z$ 至少为 $\delta$ 的紧方向集上得到 $L_K\ge1+\eta$。取充分接近一的 $\lambda<1$，使 $\lambda(1+\eta)\ge1$ 且 $(1-\lambda)R<a/4$。于是
\[
                        Q=\lambda K\cup H_\delta\subset K
\]
在所有方向上都包含完整单位线段：$Z$ 的领圈使用未收缩的选择器图像，领圈的补集则使用 $\lambda K$ 中的长弦。在每条物理射线 $\phi\in B$ 上，整个径向层 $\lambda\rho_K(\phi)<r\le\rho_K(\phi)$ 都不与 $Q$ 的两部分相交。这一层的面积为
\[
             \frac{1-\lambda^2}{2}\int_B\rho_K(\phi)^2\,d\phi>0.
\]
将有限开包络恢复的误差取为小于这个面积差，就得到面积严格更小的原始有限竞争体。当 $Z$ 为空时，一致收缩已经足够；当 $Z$ 是整个圆周时，$H$ 本身就提供所有方向的单位线段，不需要补集上的长度余量。这里比较面积的对象是这个闭选择器核；临界方向上的单侧极限单位线段也必须计入。

附录~\ref{sec:constant-length-classification} 给出此处所需的完整分类证明。引理~\ref{lem:constant-length-classification} 适用于实际边界由有限条原点同心圆弧、直段与奇异点组成的紧半代数原点星体；附录也证明每个原始有限体都满足这项边界假设。若 $L_K=1$ 在一个开方向区间上恒成立，则其中每个方向都有过原点的最长单位段。去掉有限个方向并细分后，两条相反物理射线上的半径 $r(q)=\rho_K(u_q)$、$s(q)=\rho_K(-u_q)$ 分别为常数，且 $r+s=1$。对于 $Z=\mathbb {RP}^1$ 的全圆紧约束情形，推论~\ref{cor:classification-full-circle} 进一步由 $L_*=1$ 推出有限例外外的 $L_K=1$，再用极坐标积分得到 $|K|\ge\pi/4$。附录中的更小有限竞争体由完整单位针体作面积恢复得到，因而这一个子类不能达到 $A_*$。

混合的紧约束扇区和有限个例外三角形仍可能填满可用的临界核。在这种情况下，现有论证还没有证明存在严格的剩余面积差，因此有限体是否一概不能取到下确界，仍是开放问题。带符号的端点链恒等式也不够：它的跳跃连接可能离开该体，绕数也可能对重叠部分重复计数。这些缺口具体卡在达到性的论证路线上，不能据此假定尖锐下确界有或没有有限极小体。
